\chapter{Aboriginal and Torres Strait Islander Social and Emotional Wellbeing}

\section*{Definition and Overview}
Social and emotional wellbeing is the Aboriginal and Torres Strait Islander concept of health that extends beyond the absence of illness to include connection to culture, community, country, family, and spirituality. It recognises that the wellbeing of a person is inseparable from the wellbeing of their family and community, and that the effects of colonisation, dispossession, and intergenerational trauma are carried in the present. This chapter describes the strengths of Aboriginal and Torres Strait Islander peoples, the social and emotional wellbeing framework, the mental health challenges that disproportionately affect Indigenous communities, and the culturally safe services and approaches that promote healing.

\section*{Epidemiology}
Aboriginal and Torres Strait Islander peoples experience higher rates of psychological distress, suicide, and hospitalisation for mental health conditions than non-Indigenous Australians. The suicide rate among Indigenous people is around twice that of the general population, and is highest among young people. Rates of substance use, family violence, and contact with the child protection and justice systems are also elevated, all of which reflect and compound social and emotional distress. At the same time, Indigenous peoples show remarkable resilience, cultural strength, and community solidarity, and most report good or excellent overall health. Access to culturally appropriate mental health care remains limited in many regions, particularly in remote communities.

\section*{Aetiology and Pathophysiology}
Social and emotional wellbeing is shaped by historical and ongoing social determinants. Colonisation, forced removal of children (the Stolen Generations), dispossession from land, racism, and the loss of language and culture have created intergenerational trauma that continues to affect families. Contemporary stressors include poverty, overcrowded housing, unemployment, discrimination, grief from premature deaths, and the high burden of chronic disease. These factors interact with individual experiences, family circumstances, and community strengths. Protective factors are equally powerful: strong cultural identity, connection to country, community support, language, ceremony, and the presence of Elders are associated with better wellbeing. Understanding that distress often has collective and historical roots is central to providing effective care.

\section*{Clinical Features}
\begin{itemize}
    \item Psychological distress, low mood, anxiety, and anger
    \item Grief and loss, including multiple bereavements within families and communities
    \item Substance use as a way of coping with distress
    \item Family and community conflict, and experiences of family violence
    \item Suicidal thoughts and behaviour, especially among young people
    \item Physical symptoms such as headache, fatigue, and sleep disturbance that accompany distress
    \item Loss of connection to culture, family, or country
    \item Presenting strengths: cultural knowledge, community roles, humour, and resilience
\end{itemize}

\section*{Differential Diagnosis}
Mental health assessment in this context must consider both clinical conditions and the social and emotional wellbeing framework.
\begin{itemize}
    \item \textbf{Depression and anxiety disorders:} may present with physical symptoms or irritability
    \item \textbf{Post-traumatic stress disorder:} common given high rates of trauma exposure
    \item \textbf{Complicated grief:} multiple and unresolved losses
    \item \textbf{Substance use disorders:} often intertwined with distress
    \item \textbf{Social and emotional distress:} responses to racism, discrimination, and life stressors
    \item \textbf{Physical illness:} chronic disease, pain, and sleep disorders that affect mood
\end{itemize}

\section*{When to Seek Medical Care}
People experiencing distress that is severe, persistent, or affecting daily function, relationships, or safety should seek help, ideally from a service that provides culturally safe care. Families and communities should encourage and support members to access help and can contact services on their behalf. Elders, community leaders, and Aboriginal health workers are trusted first points of contact in many communities.

\textbf{Call 000 (or local emergency services) for:}
\begin{itemize}
    \item Immediate risk of suicide or self-harm
    \item Sudden severe confusion, hallucinations, or danger to self or others
    \item Any situation in which a person's safety or the safety of others is at risk
\end{itemize}

\section*{Diagnosis and Investigations}
\begin{itemize}
    \item \textbf{Clinical assessment:} a full history with attention to cultural background, family, community, and trauma history
    \item \textbf{Culturally safe engagement:} involvement of Aboriginal health workers, interpreters, and community support
    \item \textbf{Screening tools:} validated measures of psychological distress and wellbeing, used respectfully
    \item \textbf{Risk assessment:} careful exploration of suicidal thoughts, plans, and means
    \item \textbf{Physical health review:} assessment of chronic disease, substance use, sleep, and pain
    \item \textbf{Social assessment:} housing, income, family responsibilities, and community connections
    \item \textbf{Involvement of family and community:} with the person's consent, as healing often occurs collectively
\end{itemize}

\section*{Management}
\begin{itemize}
    \item \textbf{Culturally safe care:} services delivered by or in partnership with Aboriginal and Torres Strait Islander health workers and organisations
    \item \textbf{Psychological therapies:} trauma-informed counselling and evidence-based therapies adapted to culture
    \item \textbf{Community-based healing:} programs run by Indigenous communities that strengthen identity, connection, and culture
    \item \textbf{Social and cultural activities:} connection to country, ceremony, art, sport, and language programs
    \item \textbf{Family and kinship involvement:} care that includes family members where appropriate
    \item \textbf{Medicines:} appropriate use of medication for depression, anxiety, or other conditions, with clear explanation
    \item \textbf{Integrated care:} coordination between mental health, drug and alcohol, primary care, and social services
    \item \textbf{Follow-up and continuity:} sustained relationships with the same workers and services
    \item \textbf{Crisis and suicide prevention:} community-based prevention programs and accessible crisis support
\end{itemize}

\section*{Complications}
\begin{itemize}
    \item Suicide and self-harm, with devastating effects on families and communities
    \item Substance use disorders and their health and social consequences
    \item Family violence and breakdown of relationships
    \item Involvement with the justice and child protection systems
    \item Physical health problems, including chronic disease, driven partly by stress
    \item Disengagement from education and employment
    \item Loss of cultural connection and intergenerational transmission of trauma
    \item Barriers to accessing care due to racism, distrust, and service gaps
\end{itemize}

\section*{Prognosis and Outcomes}
With culturally safe, holistic, and sustained support, significant healing is possible. Many Indigenous people recover from episodes of distress and thrive when their care respects culture, family, and community. Programs led by Aboriginal and Torres Strait Islander communities, such as social and emotional wellbeing teams, men's and women's groups, and on-country healing initiatives, show strong results. The greatest improvements follow from addressing the social determinants of health: secure housing, employment, education, and self-determination. Healing is a long-term, collective process, and hope, cultural continuity, and community strength are central to it.

\section*{Prevention and Self-Care}
\begin{itemize}
    \item Strengthen connection to culture, family, community, and country
    \item Involve Elders and community leaders in wellbeing activities
    \item Access culturally safe services early rather than waiting for crisis
    \item Take part in community programs, sport, art, and ceremony
    \item Practise self-care: sleep, physical activity, healthy food, and time on country
    \item Reduce alcohol and other drug use that may worsen distress
    \item Talk openly about distress and support others to seek help
    \item Support community-driven initiatives that address the causes of distress
    \item Health services should train staff in cultural safety and employ Indigenous workers
\end{itemize}

\section*{Sources and Further Reading}
The following reputable sources provide further information for healthcare students and patients seeking reliable, current advice about Aboriginal and Torres Strait Islander social and emotional wellbeing.
\begin{itemize}
    \item HealthInfoNet. \textit{Social and emotional wellbeing resources.} https://healthinfonet.ecu.edu.au/
    \item Australian Indigenous HealthInfoNet. \textit{Wellbeing and mental health.} https://healthinfonet.ecu.edu.au/
    \item Gayaa Dhuwi (Proud Spirit) Australia. \textit{Indigenous mental health leadership.} https://gayaadhuwi.org.au/
    \item Head to Health. \textit{Support services.} https://www.headtohealth.gov.au/
    \item 13YARN. \textit{Crisis support line for Aboriginal and Torres Strait Islander people.} https://www.13yarn.com.au/
    \item Department of Health and Aged Care. \textit{Social and emotional wellbeing programs.} https://www.health.gov.au/
\end{itemize}
